\problemname{Skiing}
Johan likes skiing.
By that we do not mean downhill skiing, which Johan is in fact very afraid of.
Cross-country skiing, however, that is his thing.
When cross-country skiing, one needs large flat areas.

Johan has surveyed a large rectangular area out in the forest, which has rather uneven terrain.
Here, Johan wants to select a certain square to ski around on, that is large enough to make the skiing interesting.
The square shall have exactly size $L \times L$, and have sides parallel to the area's sides.

Now he asks you to find such a square.
For it to be suitable for cross-country skiing, he has two requirements.
First, there must be no trees in the square, and second, the height difference between the highest and lowest point in this square shall be as small as possible.

If there are multiple possible such squares, you shall first choose the one that is furthest north, i.e.\ has the lowest row number. If there are still multiple possibilities,
you shall secondly choose the one that is furthest west, i.e.\ has the lowest column number.

\section*{Input}
The first line contains three integers $R$, $C$, $L$ such that $1\leq R,C \leq 1000, 1 \leq L \leq \min(R,C)$.
$R$ is the number of rows in the large area, $C$ the number of columns, and $L$ the size of the desired square.

Then follow $R$ lines, one for each row in the area. A line contains $C$ integers, one for each column in the area.

The $c$th number on the $r$th line describes the height $H_{rc}$ at that point in the area, such that $-1 \leq H_{rc} \leq 10^9$.
If $H_{rc} = -1$, it means there is a \textbf{tree at that location} instead.

\section*{Output}
Find $r_l$, $c_l$ such that Johan's square spans the coordinates $r_l \leq r <r_l + L$, $c_l \leq c < c_l + L$. $r_l$ and $c_l$ shall be 0-indexed; for example, $r_l = 0$ means the first row (furthest north), and $c_l = 0$ means the first column (furthest west).

It is guaranteed that a solution exists.

\section*{Scoring}
Your solution will be tested on a set of test groups. To earn points for a group, you must pass all test cases in that group.

\begin{tabular}{|l|l|l|}
\hline
Group & Points & Constraints \\ \hline
 1     & 21         &  $0 \leq H_{rc} \leq 1$ \\ \hline
 2     & 36         &  $R,C \leq 50$ \\ \hline
 3     & 27         &  $R,C \leq 300$ \\ \hline
 4     & 16         &  $R,C \leq 1000$ \\ \hline
\end{tabular}
